\documentclass{article}
\usepackage{graphicx} % Required for inserting images
\usepackage{amsmath}
\usepackage{cancel}

\title{Analysis PS2}
\author{Anthony, Jordan, Kemin}
\date{August 2026}

\begin{document}

\maketitle

\section*{Q1}

\subsection*{Case 1: $x\ge 0, y \ge 0$}
$||x|-|y|\le |x-y|$\\
$|x-y|\le |x-y|$
\subsection*{Case 2: $x<0, y<0$}
$||x|-|y|\le |x-y||$\\
$|y-x|\le |x-y|$\\
$|-(x-y)|\le |x-y|$\\
$|x - y|\le |x-y|$
\subsection*{Case 3: $x\ge 0, y<0$}
$||x|-|y|\le |x-y||$\\
$|x+y| \le |x-y|$
\subsection*{Case 4: $x<0, y\ge 0$}
$||x|-|y|\le |x-y||$\\
$|-x-y| \le |x-y|$
\section*{Q2}
$T_a(x) = e^a+e^a(x-a)+\frac{e^a}{2!}(x-a)^2+\frac{e^a}{3!}(x-a)^3+\frac{e^a}{4!}(x-a)^4+\frac{e^a}{5!}(x-a)^5$\\
$T_0(x)=1+x+\frac{x^2}{2!}+\frac{x^3}{3!}+\frac{x^4}{4!}+\frac{x^5}{5!}$\\
$T_0(1)=1+1+\frac{1}{2!}+\frac{1}{3!}+\frac{1}{4!}+\frac{1}{5!}$\\
$T_0(1)=\frac{120+120+60+20+5+1}{120}$\\
$T_0(1)=\frac{326}{120}$

\section*{Q3}
$\frac{d\log (x(t))}{dt}$\\
$\frac{1}{x}\frac{dx(t)}{dt}$\\
$\frac{}{}$

\section*{Q4}

\subsection*{4.1}

\[
H_f=
\begin{bmatrix}
-\dfrac{1}{4}\,x^{-3/2} & 0\\[10pt]
0 & -\dfrac{1}{4}\,y^{-3/2}
\end{bmatrix}
=\frac{1}{4}
\begin{bmatrix}
-\dfrac{1}{x^{3/2}} & 0\\[10pt]
0 & -\dfrac{1}{y^{3/2}}
\end{bmatrix}
\]

\[
|Hf|=\frac{1}{16}\cdot\frac{1}{(xy)^{3/2}}>0
\]

The Hessian is negative definite, so $f$ is strictly concave over
$\mathbb{R}^2_{++}$, its domain.

\subsection*{4.2 \quad $g(x,y)=\sqrt{xy}$}

\[
Hg=
\begin{bmatrix}
-\dfrac{\sqrt{y}}{4}\,x^{-3/2} & \dfrac{1}{4\sqrt{xy}}\\[10pt]
\dfrac{1}{4\sqrt{xy}} & -\dfrac{\sqrt{x}}{4}\,y^{-3/2}
\end{bmatrix}
=\frac{1}{4}
\begin{bmatrix}
-\dfrac{y^{1/2}}{x^{3/2}} & \dfrac{1}{\sqrt{xy}}\\[10pt]
\dfrac{1}{\sqrt{xy}} & -\dfrac{x^{1/2}}{y^{3/2}}
\end{bmatrix}
\]

\[
|Hg|=\frac{1}{xy}-\frac{1}{xy}=0
\]

The Hessian is negative semidefinite, so $g$ is concave but not strictly concave over
$\mathbb{R}^2_{++}$, its domain.

\section*{Q5}

\subsection*{5.1}

\[
\bigcap_{n=1}^{\infty} I_n=\{0\},
\]
because for any $z\in\mathbb{R}$, $z\neq 0$, there exists $n\in\mathbb{N}$ such
that
\[
d(0,\tfrac{1}{n})<d(0,z)
\quad\Longleftrightarrow\quad
\left|\tfrac{1}{n}\right|<|z|
\quad\Longrightarrow\quad
z\notin I_n .
\]
And $0\in I_n$ because $0<\tfrac{1}{n}$ for every $n\in\mathbb{N}$.

$\{0\}$ is a closed set because it is the complement of
$(-\infty,0)\cup(0,\infty)$, which is the union of two open balls and thus an
open set by way of Theorems 1 and 2.

\subsection*{5.2}

For center $x=\tfrac{1}{2}$ and radius $r=\dfrac{n-1}{n+1}$,
\[
\left(\tfrac{1}{2}-\tfrac{1}{n+1}=\tfrac{n-1}{n+1}\right)
\]
we have $\forall y\in J_n$, $d(x,y)\le r$. Also, the complement of $J_n$ is the
open set
\[
\left(-\infty,\tfrac{1}{n+1}\right)\cup\left(\tfrac{n}{n+1},\infty\right).
\]
Taken together, that implies that $J_n$ is a closed ball.

For $\displaystyle\bigcup_{i=1}^{\infty}J_n$:

$0<\dfrac{1}{n+1}$ for any $n$, so for any $u\le 0$,
$u\notin\displaystyle\bigcup_{i=1}^{\infty}J_n$.

Similarly $1>1-\dfrac{1}{n+1}$ for any $n$, so for any $v\ge 1$,
$v\notin\displaystyle\bigcup_{i=1}^{\infty}J_n$.

\section*{Q6}

\section*{Q7}

Bordered Hessian:
\[
H=
\begin{bmatrix}
0 & \alpha AK^{\alpha-1}L^{1-\alpha} & (1-\alpha)AK^{\alpha}L^{-\alpha}\\[4pt]
\alpha AK^{\alpha-1}L^{1-\alpha} & (\alpha^2-\alpha)AK^{\alpha-2}L^{1-\alpha} & (\alpha-\alpha^2)AK^{\alpha-1}L^{-\alpha}\\[4pt]
(1-\alpha)AK^{\alpha}L^{-\alpha} & (\alpha-\alpha^2)AK^{\alpha-1}L^{-\alpha} & (\alpha^2-\alpha)AK^{\alpha}L^{-\alpha-1}
\end{bmatrix}
\]
where
\[
a=\alpha AK^{\alpha-1}L^{1-\alpha},\qquad
b=(1-\alpha)AK^{\alpha}L^{-\alpha},\qquad
c=(\alpha^2-\alpha)AK^{\alpha-2}L^{1-\alpha},
\]
\[
d=(\alpha-\alpha^2)AK^{\alpha-1}L^{-\alpha},\qquad
e=(\alpha^2-\alpha)AK^{\alpha}L^{-\alpha-1}.
\]

\subsection*{Leading principal minors}

\[
|H_2|=-\alpha^2A^2K^{2\alpha-2}L^{2-2\alpha}\le 0
\qquad\text{for }K,L\ge 0
\]

\[
|H_3|=2abd-a^2e-b^2c
\]
\[
=2(\alpha)(1-\alpha)(\alpha-\alpha^2)A^3K^{3\alpha-2}L^{1-3\alpha}
-(\alpha^4-\alpha^3)A^3K^{3\alpha-2}L^{1-3\alpha}
\]
\[
\qquad-(1-\alpha)^2(\alpha^2-\alpha)A^3K^{3\alpha-2}L^{1-3\alpha},
\]
need to show $\ge 0$ for $K,L\ge 0$.

Divide by $(1-\alpha)A^3$:
\[
2\alpha^2(1-\alpha)\cancel{K^{3\alpha-2}L^{1-3\alpha}}
\ \ge\ -\alpha^3\cancel{K^{3\alpha-2}L^{1-3\alpha}}
+(1-\alpha)(\alpha^2-\alpha)\cancel{K^{3\alpha-2}L^{1-3\alpha}}
\]
\[
2\alpha^2\ \ge\ (1-\alpha)(\alpha^2-\alpha)-\alpha^3
\]
\[
2\alpha\ \ge\ (1-\alpha)(\alpha-1)-\alpha^2
\]
\[
2\alpha\ \ge\ -\big((\alpha-1)^2+\alpha^2\big)
\]
\[
2\alpha\ \ge\ -\big(2\alpha^2-2\alpha+1\big)
\]
\[
0\ \ge\ -2\alpha^2-1\quad\checkmark
\]


\section*{Q8}



\section*{Q9}

\begin{enumerate}
\item $A$ is bounded, $M=\sqrt{10}$, $\bar{x}=0$.

\item $B$ is not bounded. All $x<0$ are in $B$.

\item $C$ is not bounded. Assuming $\mathbb{R}_+$ means non-negative, all
$(0,y)$ and $(x,0)$ are in $C$.

\item $D$ is bounded, $M=10$ and $\bar{x}=(0,0)$.
\end{enumerate}


\end{document}
